-\sum_i\gamma_iB_0m_i
+
\sum_iQ_i\!\left(m_i^2-\tfrac23\right)
\right]
\\
&\quad
-
\left[
-\sum_i\omega_i m_i
+
\sum_iQ_i\!\left(m_i^2-\tfrac23\right)
\right]
\\
&=
\Lambda_s
-
\sum_i m_i(\gamma_iB_0-\omega_i).
\end{aligned}
\label{eq:Lambda_config_general}
\end{equation}
this becomes
\begin{equation}
\begin{aligned}
&e^{iH_{\mathrm{free}}t}
\sigma_x
e^{-iH_{\mathrm{free}}t}
\notag\\
&=
\sum_{\mathbf m}
\left[
\sigma_+ e^{i\Lambda(\mathbf m)t}
+
\sigma_- e^{-i\Lambda(\mathbf m)t}
\right]
\otimes
\ket{\mathbf m}\!\bra{\mathbf m}
\\
&=
\sum_{\mathbf m}
\big[
\sigma_x\cos(\Lambda(\mathbf m)t)
-
\sigma_y\sin(\Lambda(\mathbf m)t)
\big]
\otimes
\ket{\mathbf m}\!\bra{\mathbf m}.
\end{aligned}
\label{eq:sigmax_IP_step}
\end{equation}



Consequently, the bare electronic part of the microwave interaction becomes
\begin{equation}
\begin{aligned}
H_{\mathrm{MW}}^{(I)}(t)
&=
\frac{\gamma_e B_{1,\mathrm{MW}}(t)}{\sqrt2}
\sum_{\mathbf m}
\big[
\sigma_x\cos(\Lambda(\mathbf m)t)
-
\sigma_y\sin(\Lambda(\mathbf m)t)
\big]
\\
&\hspace{1cm}\otimes
\ket{\mathbf m}\!\bra{\mathbf m}.
\end{aligned}
\label{eq:HMW_I}
\end{equation}
Here, the nuclear configuration is preserved to leading order during the bare
microwave-driven electronic transition.\\

%



We now transform the second term in the dressed microwave interaction,
\[
H_{\mathrm{MW},eN}
=
-\frac{\gamma_e B_{1,\mathrm{MW}}(t)}{\sqrt2}
\sigma_y\otimes G .
\]
Its interaction-picture form is
\begin{equation}
H_{\mathrm{MW},eN}^{(I)}(t)
=
-\frac{\gamma_e B_{1,\mathrm{MW}}(t)}{\sqrt2}
e^{iH_{\mathrm{free}}t}
\left(
\sigma_y\otimes G
\right)
e^{-iH_{\mathrm{free}}t}.
\label{eq:HeN_IP_start}
\end{equation}

Since the electron-nuclear modulation term contributes only at first order in the small misalignment angles \(\theta_i\), we employ a leading-order description of its interaction-picture evolution. For the electronic part, we retain the dominant bare precession at \(\Lambda_s\) rather than the full configuration-dependent transition frequency \(\Lambda(\mathbf m)\). For the nuclear part, we approximate the manifold-dependent precession by the effective mean generator
\[
H_N
=
-\sum_i \delta_i I_{iz},
\,\, \text{with}\,\,
\delta_i
=
\frac{\gamma_iB_0+\omega_i}{2}.
\]
This captures the leading averaged dynamics of the weak modulation without
tracking the full configuration-resolved evolution.

For the same reason, we also approximate the transformation as
\begin{equation}
\begin{aligned}
&e^{iH_{\mathrm{free}}t}
\left(
\sigma_y\otimes G
\right)
e^{-iH_{\mathrm{free}}t}
\notag\\
&\approx
\left(
e^{i\Lambda_s\sigma_z t/2}
\sigma_y
e^{-i\Lambda_s\sigma_z t/2}
\right)
\otimes
\left(
e^{iH_N t}
G
e^{-iH_N t}
\right).
\end{aligned}
\label{eq:sigmayG_factorized}
\end{equation}

The electronic part evaluates to
\[
e^{i\Lambda_s\sigma_z t/2}
\sigma_y
e^{-i\Lambda_s\sigma_z t/2}
=
\sigma_x\sin(\Lambda_s t)
+
\sigma_y\cos(\Lambda_s t),
\]
while the nuclear operator
\[
G
=
\sum_i\theta_i
\left(
-I_{ix}\sin\phi_i
+
I_{iy}\cos\phi_i
\right)
\]
evolves under \(H_N\) as
\[
e^{iH_N t}
G
e^{-iH_N t}
=
\sum_i\theta_i
\left[
I_{ix}\cos(\delta_i t-\phi_i)
+
I_{iy}\sin(\delta_i t-\phi_i)
\right].
\]

Therefore, to leading order,
\begin{equation}
\begin{aligned}
H_{\mathrm{MW},eN}^{(I)}(t)
&\approx
-\frac{\gamma_e B_{1,\mathrm{MW}}(t)}{\sqrt2}
\big[
\sigma_x\sin(\Lambda_s t)
+
\sigma_y\cos(\Lambda_s t)
\big]
\\
&\quad\otimes
\sum_i\theta_i
\big[
I_{ix}\cos(\delta_i t-\phi_i)
+
I_{iy}\sin(\delta_i t-\phi_i)
\big].
\end{aligned}
\label{eq:HeN_I}
\end{equation}


\paragraph{Rotating-wave approximation (RWA).}
We drive the electron at microwave frequency $\omega_{\mathrm{MW}}$ with a magnetic field
\begin{equation}
B_{1,\mathrm{MW}}(t)
=
\frac{\sqrt2\,\Omega(t)}{\gamma_e}
\cos(\omega_{\mathrm{MW}} t),
\end{equation}
where $\Omega(t)$ denotes the effective time-dependent Rabi-frequency envelope. 
Substituting this field into Eq.~\eqref{eq:HMW_I} gives the prefactor
\[
\frac{\gamma_e B_{1,\mathrm{MW}}(t)}{\sqrt2}
=
\Omega(t)\cos(\omega_{\mathrm{MW}} t).
\]
Using the product-to-sum identity $\cos\alpha\cos\beta=\tfrac12[\cos(\alpha-\beta)+\cos(\alpha+\beta)]$
and $\cos\alpha\sin\beta=\tfrac12[\sin(\beta+\alpha)+\sin(\beta-\alpha)]$,
the electron interaction-picture drive
\(
\propto \sigma_x\cos[\Lambda(\mathbf m)t]-\sigma_y\sin[\Lambda(\mathbf m)t]
\)
splits into a slow (difference-frequency) piece and explicit counter-rotating (sum-frequency) pieces:
\begin{align}
&\notag\hspace{-0.0cm}H_{e}^{(\mathrm{full})}(t)= \frac{\Omega(t)}{2}\sum_{\mathbf m}\\
&\hspace{-0.2cm}
\Big\{\underbrace{\big[\sigma_x\cos((\omega_{\mathrm{MW}}-\Lambda(\mathbf m))\,t)-\sigma_y\sin((\omega_{\mathrm{MW}}-\Lambda(\mathbf m))\,t)\big]}_{\text{slow (near-resonant)}}\nonumber\\
&\hspace{-0.2cm}
+\notag \underbrace{\big[\sigma_x\cos\big((\omega_{\mathrm{MW}}+\Lambda(\mathbf m))t\big)- 
\hspace{0.2cm} \sigma_y\sin\big((\omega_{\mathrm{MW}}+\Lambda(\mathbf m))t\big)\big]}_{\text{fast (counter-rotating)}}\Big\}\\
&\otimes \ket{\mathbf m}\!\bra{\mathbf m}.
\label{eq:He_full_split}
\end{align}
The \emph{rotating-wave approximation} (RWA) keeps only the slow terms and discards the counter-rotating term at  $(\omega_{\mathrm{MW}}+\Lambda(\mathbf m))t\big)$, yielding
\begin{equation}
\boxed{
\begin{aligned}
&H_{e}^{(\mathrm{RWA})}(t)
= \frac{\Omega(t)}{2} \sum_{\mathbf m}
\big[\sigma_x\cos((\omega_{\mathrm{MW}}-\Lambda(\mathbf m))\,t)\\ & \hspace{1.9cm}-\sigma_y\sin((\omega_{\mathrm{MW}}-\Lambda(\mathbf m))\,t)\big]
\otimes \ket{\mathbf m}\!\bra{\mathbf m}.
\end{aligned}
}
\label{eq:HeRWA_time_nophase_Lambda_m}
\end{equation}
Applying the same procedure to the electron–nuclear modulation term Eq.~\eqref{eq:HeN_I}, gives
\begin{align}
&\notag H_{eN}^{(\mathrm{full})}(t)\\
&\notag=
\underbrace{-\,\frac{\Omega(t)}{2}\,
\big[\sigma_x\sin((\omega_{\mathrm{MW}}-\Lambda_s) t)+\sigma_y\cos((\omega_{\mathrm{MW}}-\Lambda_s) t)\big]}_{\text{slow (near-resonant in the electron sector)}}\;\\
&\notag\;\otimes\;
\sum_i \theta_i\big[I'_{ix}\cos(\delta_i t-\phi_i)+I'_{iy}\sin(\delta_i t-\phi_i)\big]\\
&\notag\hspace{-0.0cm}\quad
\underbrace{-\,\frac{\Omega(t)}{2}\,
\big[\sigma_x\sin\big((\Lambda_s+\omega_{\mathrm{MW}})t\big)
      +\sigma_y\cos\big((\Lambda_s+\omega_{\mathrm{MW}})t\big)\big]}_{\text{fast (counter-rotating)}}\\
&\;\otimes\;
\sum_i \theta_i\big[I'_{ix}\cos(\delta_i t-\phi_i)+I'_{iy}\sin(\delta_i t-\phi_i)\big].
\label{eq:HeN_full_split}
\end{align}
Discarding the fast rotating term again, gives 
\begin{equation}
\boxed{
\begin{aligned}
&H_{eN}^{(\mathrm{RWA})}(t)\\
&\notag=-\,\frac{\Omega(t)}{2}\,
\big[\sigma_x\sin((\omega_{\mathrm{MW}}-\Lambda_s) t)+\sigma_y\cos((\omega_{\mathrm{MW}}-\Lambda_s) t)\big]\\
&\;\otimes\;
\sum_i \theta_i\big[I'_{ix}\cos(\delta_i t-\phi_i)+I'_{iy}\sin(\delta_i t-\phi_i)\big].
\end{aligned}
}
\label{eq:HeN_RWA_time_nophase_final}
\end{equation}

\section{NV-register simulation parameters}
\label{app:nv_params}
\noindent
The NV electronic spin ($S{=}1$), with gyromagnetic ratio 
$\gamma_e/2\pi = 28.024~\text{GHz/T}$, couples to one intrinsic 
$^{14}$N nucleus and two proximate $^{13}$C spins through anisotropic 
hyperfine tensors $A_{zz}$ and $A_{\perp}$, with an external magnetic 
field $B_0 = 0.45~\text{T}$ aligned with the NV axis. 
The nuclear-spin parameters listed in Table~\ref{tab:nv_params_table} 
correspond to typical experimental ranges reported in 
Refs.~\cite{Chen,nizovtsev}.\\
\begin{table}[t]
\centering
\caption{Physical parameters used for the nuclear spins in the 
NV-$^{14}$N-$^{13}$C$_1$-$^{13}$C$_2$ register simulation.}
\label{tab:nv_params_table}
\begin{tabular}{lccc}
\toprule
Parameter & $^{14}$N ($I{=}1$) & $^{13}$C$_1$ ($I{=}\tfrac{1}{2}$) & $^{13}$C$_2$ ($I{=}\tfrac{1}{2}$) \\
\midrule
$\gamma/2\pi$ (MHz/T) & 3.077 & 10.71 & 10.71 \\
$A_{zz}/2\pi$ (MHz)                     & $-2.14$ & 2.281 & $-1.011$ \\
$A_{\perp}/2\pi$ (MHz)                  & 0.00 & 0.240 & 0.014 \\
$Q/2\pi$ (MHz)                          & $-5.01$ & 0.00 & 0.00 \\
\bottomrule
\end{tabular}
\end{table}


